% GENERATED FILE. Edit canonical lessons, then run npm run generate:books.
% canonical-source-hash: fnv1a64:41dedccf24559cbe
% canonical-lessons: KA-C387-savalu, KA-C387-misalidu, KA-C387-samatolana, KA-C387-adyate, KA-C387-maulya

\chapter{Priorities and Values}
\label{ch:ka-priorities-and-values}
\hlchaptermodality{387}

\begin{chapteropening}
\textbf{By the end of this chapter:} \emph{I can say a challenge, to set aside, balance, a priority and a value.}

Complete the last lesson of chapter 387: I can say a challenge, to set aside, balance, a priority and a value.
\end{chapteropening}

\section[\texorpdfstring{\kn{ಸವಾಲು} (savālu)}{ಸವಾಲು (savālu)}]{\kn{ಸವಾಲು} (savālu) — a challenge}

\label{lesson:KA-C387-savalu}



\begin{quote}
Before the new one: say the Kannada for to portray, then the Kannada for blame.
\end{quote}

\subsection*{You'll want to know: \kn{ಸವಾಲು}}
\textbf{\kn{ಸವಾಲು}} — \emph{savālu} — \textquotedblleft{}a challenge\textquotedblright{}.

\begin{quote}
\kn{ಹೊಸ} \kn{ಊರಿನಲ್ಲಿ} \kn{ಕೆಲಸ} \kn{ಹುಡುಕುವುದು} \kn{ದೊಡ್ಡ} \kn{ಸವಾಲು}.
\end{quote}

\begin{quote}
\emph{hosa ūrinalli kelasa huḍukuvudu doḍḍa savālu.}
\end{quote}

\begin{quote}
Finding work in a new town is a big challenge.
\end{quote}

\kn{ಸವಾಲು} \kn{ಹಾಕು}, to throw down a challenge.

\begin{cousinweb}
\textbf{\kn{ಸವಾಲು}} (\emph{savālu}) is a loan from Hindustani \emph{savāl}, which comes from Arabic \emph{suʾāl}, \textquotedblleft{}a question\textquotedblright{}; Kannada uses it for a challenge.
\end{cousinweb}

\begin{grammarlens}[title={What you've built}]
One more word for giving reasons and opinions.
\end{grammarlens}

\subsection*{Your turn}
\begin{itemize}
\raggedright
  \item \emph{Say it:} \emph{savālu}
  \item \emph{Say it:} \emph{savālu}, once more
  \item \emph{Say it:} \emph{hosa ūrinalli kelasa huḍukuvudu doḍḍa savālu.} — Finding work in a new town is a big challenge.
  \item \emph{Recall:} say the Kannada for to portray, then the Kannada for blame, then say \emph{savālu} again
\end{itemize}

\begin{tcolorbox}[breakable,colback=teal!4,colframe=teal!35!black,title={Before you move on}]
What does \kn{ಸವಾಲು} mean? (\textquotedblleft{}A challenge\textquotedblright{}.) Say it once more.
\end{tcolorbox}
\section[\texorpdfstring{\kn{ಮೀಸಲಿಡು} (mīsaliḍu)}{ಮೀಸಲಿಡು (mīsaliḍu)}]{\kn{ಮೀಸಲಿಡು} (mīsaliḍu) — to set aside}

\label{lesson:KA-C387-misalidu}



\begin{quote}
Before the new one: say the Kannada for blame, then the Kannada for a challenge.
\end{quote}

\subsection*{You'll want to know: \kn{ಮೀಸಲಿಡು}}
\textbf{\kn{ಮೀಸಲಿಡು}} — \emph{mīsaliḍu} — \textquotedblleft{}to set aside\textquotedblright{}.

\begin{quote}
\kn{ನಾನು} \kn{ಪ್ರತಿ} \kn{ತಿಂಗಳು} \kn{ಸ್ವಲ್ಪ} \kn{ಹಣ} \kn{ಮೀಸಲಿಟ್ಟೆ}.
\end{quote}

\begin{quote}
\emph{nānu prati tiṅgaḷu svalpa haṇa mīsaliṭṭe.}
\end{quote}

\begin{quote}
I set aside a little money every month.
\end{quote}

Past (I …): \kn{ಮೀಸಲಿಟ್ಟೆ}, \emph{mīsaliṭṭe}.

\kn{ಮೀಸಲು} (kept apart) + \kn{ಇಡು} (to put).

\begin{cousinweb}
\textbf{\kn{ಮೀಸಲಿಡು}} (\emph{mīsaliḍu}) is native Dravidian, built from \emph{mīsal}, \textquotedblleft{}something set apart and left untouched\textquotedblright{} (DEDR 4841), and \emph{iḍu}, \textquotedblleft{}to put\textquotedblright{} (DEDR 442).
\end{cousinweb}

\begin{grammarlens}[title={What you've built}]
One more word for giving reasons and opinions.
\end{grammarlens}

\subsection*{Your turn}
\begin{itemize}
\raggedright
  \item \emph{Say it:} \emph{mīsaliḍu}
  \item \emph{Say it:} \emph{mīsaliḍu}, once more
  \item \emph{Say it:} \emph{nānu prati tiṅgaḷu svalpa haṇa mīsaliṭṭe.} — I set aside a little money every month.
  \item \emph{Recall:} say the Kannada for blame, then the Kannada for a challenge, then say \emph{mīsaliḍu} again
\end{itemize}

\begin{tcolorbox}[breakable,colback=teal!4,colframe=teal!35!black,title={Before you move on}]
What does \kn{ಮೀಸಲಿಡು} mean? (\textquotedblleft{}To set aside\textquotedblright{}.) Say it once more.
\end{tcolorbox}
\section[\texorpdfstring{\kn{ಸಮತೋಲನ} (samatōlana)}{ಸಮತೋಲನ (samatōlana)}]{\kn{ಸಮತೋಲನ} (samatōlana) — balance}

\label{lesson:KA-C387-samatolana}



\begin{quote}
Before the new one: say the Kannada for a challenge, then the Kannada for to set aside.
\end{quote}

\subsection*{You'll want to know: \kn{ಸಮತೋಲನ}}
\textbf{\kn{ಸಮತೋಲನ}} — \emph{samatōlana} — \textquotedblleft{}balance\textquotedblright{}.

\begin{quote}
\kn{ಕೆಲಸ} \kn{ಮತ್ತು} \kn{ಮನೆಯ} \kn{ನಡುವೆ} \kn{ಸಮತೋಲನ} \kn{ಬೇಕು}.
\end{quote}

\begin{quote}
\emph{kelasa mattu maneya naḍuve samatōlana bēku.}
\end{quote}

\begin{quote}
We need a balance between work and home.
\end{quote}

\kn{ಸಮ} (even) + \kn{ತೋಲನ} (weighing).

\begin{cousinweb}
\textbf{\kn{ಸಮತೋಲನ}} (\emph{samatōlana}) is built from Sanskrit \emph{sama}, \textquotedblleft{}even, equal\textquotedblright{}, and \emph{tolana}, \textquotedblleft{}weighing\textquotedblright{} (Monier-Williams).
\end{cousinweb}

\begin{grammarlens}[title={What you've built}]
One more word for giving reasons and opinions.
\end{grammarlens}

\subsection*{Your turn}
\begin{itemize}
\raggedright
  \item \emph{Say it:} \emph{samatōlana}
  \item \emph{Say it:} \emph{samatōlana}, once more
  \item \emph{Say it:} \emph{kelasa mattu maneya naḍuve samatōlana bēku.} — We need a balance between work and home.
  \item \emph{Recall:} say the Kannada for a challenge, then the Kannada for to set aside, then say \emph{samatōlana} again
\end{itemize}

\begin{tcolorbox}[breakable,colback=teal!4,colframe=teal!35!black,title={Before you move on}]
What does \kn{ಸಮತೋಲನ} mean? (\textquotedblleft{}Balance\textquotedblright{}.) Say it once more.
\end{tcolorbox}
\section[\texorpdfstring{\kn{ಆದ್ಯತೆ} (ādyate)}{ಆದ್ಯತೆ (ādyate)}]{\kn{ಆದ್ಯತೆ} (ādyate) — a priority}

\label{lesson:KA-C387-adyate}



\begin{quote}
Before the new one: say the Kannada for to set aside, then the Kannada for balance.
\end{quote}

\subsection*{You'll want to know: \kn{ಆದ್ಯತೆ}}
\textbf{\kn{ಆದ್ಯತೆ}} — \emph{ādyate} — \textquotedblleft{}a priority\textquotedblright{}.

\begin{quote}
\kn{ನನ್ನ} \kn{ಮೊದಲ} \kn{ಆದ್ಯತೆ} \kn{ಮಕ್ಕಳ} \kn{ಓದು}.
\end{quote}

\begin{quote}
\emph{nanna modala ādyate makkaḷa ōdu.}
\end{quote}

\begin{quote}
My first priority is the children's studies.
\end{quote}

\kn{ಆದ್ಯತೆ} \kn{ಕೊಡು}, to give priority.

\begin{cousinweb}
\textbf{\kn{ಆದ್ಯತೆ}} (\emph{ādyate}) is built on Sanskrit \emph{ādya}, \textquotedblleft{}being at the beginning, first\textquotedblright{} (Monier-Williams), with -\emph{tā}, \textquotedblleft{}-ness\textquotedblright{}: firstness.
\end{cousinweb}

\begin{grammarlens}[title={What you've built}]
One more word for giving reasons and opinions.
\end{grammarlens}

\subsection*{Your turn}
\begin{itemize}
\raggedright
  \item \emph{Say it:} \emph{ādyate}
  \item \emph{Say it:} \emph{ādyate}, once more
  \item \emph{Say it:} \emph{nanna modala ādyate makkaḷa ōdu.} — My first priority is the children's studies.
  \item \emph{Recall:} say the Kannada for to set aside, then the Kannada for balance, then say \emph{ādyate} again
\end{itemize}

\begin{tcolorbox}[breakable,colback=teal!4,colframe=teal!35!black,title={Before you move on}]
What does \kn{ಆದ್ಯತೆ} mean? (\textquotedblleft{}A priority\textquotedblright{}.) Say it once more.
\end{tcolorbox}
\section[\texorpdfstring{\kn{ಮೌಲ್ಯ} (maulya)}{ಮೌಲ್ಯ (maulya)}]{\kn{ಮೌಲ್ಯ} (maulya) — a value; worth}

\label{lesson:KA-C387-maulya}



\begin{quote}
Before the new one: say the Kannada for balance, then the Kannada for a priority.
\end{quote}

\subsection*{You'll want to know: \kn{ಮೌಲ್ಯ}}
\textbf{\kn{ಮೌಲ್ಯ}} — \emph{maulya} — \textquotedblleft{}a value; worth\textquotedblright{}.

\begin{quote}
\kn{ಹಳೆಯ} \kn{ನಾಣ್ಯಗಳಿಗೆ} \kn{ಈಗ} \kn{ಹೆಚ್ಚು} \kn{ಮೌಲ್ಯ} \kn{ಇದೆ}.
\end{quote}

\begin{quote}
\emph{haḷeya nāṇyagaḷige īga heccu maulya ide.}
\end{quote}

\begin{quote}
Old coins are worth more now.
\end{quote}

\kn{ಮೌಲ್ಯಗಳು}, values (moral or social).

\begin{cousinweb}
\textbf{\kn{ಮೌಲ್ಯ}} (\emph{maulya}) is borrowed from Sanskrit \emph{maulya}, \textquotedblleft{}price\textquotedblright{} (CDIAL 10373), a form of \emph{mūlya}, \textquotedblleft{}worth\textquotedblright{}.
\end{cousinweb}

\begin{grammarlens}[title={What you've built}]
One more word for giving reasons and opinions.
\end{grammarlens}

\subsection*{Your turn}
\begin{itemize}
\raggedright
  \item \emph{Say it:} \emph{maulya}
  \item \emph{Say it:} \emph{maulya}, once more
  \item \emph{Say it:} \emph{haḷeya nāṇyagaḷige īga heccu maulya ide.} — Old coins are worth more now.
  \item \emph{Recall:} say the Kannada for balance, then the Kannada for a priority, then say \emph{maulya} again
\end{itemize}

\begin{tcolorbox}[breakable,colback=teal!4,colframe=teal!35!black,title={Before you move on}]
What does \kn{ಮೌಲ್ಯ} mean? (\textquotedblleft{}A value; worth\textquotedblright{}.) Say it once more.
\end{tcolorbox}
